\documentclass[a4paper,12pt,twoside]{memoir}

% Castellano
\usepackage[spanish,es-tabla]{babel}
\selectlanguage{spanish}
\usepackage[utf8]{inputenc}
\usepackage[T1]{fontenc}
\usepackage{lmodern} % scalable font
\usepackage{microtype}
\usepackage{placeins}
\usepackage{lscape}


\RequirePackage{booktabs}
\RequirePackage[table]{xcolor}
\RequirePackage{xtab}
\RequirePackage{multirow}

% Links
\PassOptionsToPackage{hyphens}{url}\usepackage[colorlinks]{hyperref}
\hypersetup{
	allcolors = {red}
}

% Ecuaciones
\usepackage{amsmath}

% Rutas de fichero / paquete
\newcommand{\ruta}[1]{{\sffamily #1}}

% Párrafos
\nonzeroparskip

% Huérfanas y viudas
\widowpenalty100000
\clubpenalty100000

% Evitar solapes en el header
\nouppercaseheads

% Imagenes
\usepackage{graphicx}
\newcommand{\imagen}[2]{
	\begin{figure}[!h]
		\centering
		\includegraphics[width=0.9\textwidth]{#1}
		\caption{#2}\label{fig:#1}
	\end{figure}
	\FloatBarrier
}

\newcommand{\imagenflotante}[2]{
	\begin{figure}%[!h]
		\centering
		\includegraphics[width=0.9\textwidth]{#1}
		\caption{#2}\label{fig:#1}
	\end{figure}
}



% El comando \figura nos permite insertar figuras comodamente, y utilizando
% siempre el mismo formato. Los parametros son:
% 1 -> Porcentaje del ancho de página que ocupará la figura (de 0 a 1)
% 2 --> Fichero de la imagen
% 3 --> Texto a pie de imagen
% 4 --> Etiqueta (label) para referencias
% 5 --> Opciones que queramos pasarle al \includegraphics
% 6 --> Opciones de posicionamiento a pasarle a \begin{figure}
\newcommand{\figuraConPosicion}[6]{%
  \setlength{\anchoFloat}{#1\textwidth}%
  \addtolength{\anchoFloat}{-4\fboxsep}%
  \setlength{\anchoFigura}{\anchoFloat}%
  \begin{figure}[#6]
    \begin{center}%
      \Ovalbox{%
        \begin{minipage}{\anchoFloat}%
          \begin{center}%
            \includegraphics[width=\anchoFigura,#5]{#2}%
            \caption{#3}%
            \label{#4}%
          \end{center}%
        \end{minipage}
      }%
    \end{center}%
  \end{figure}%
}

%
% Comando para incluir imágenes en formato apaisado (sin marco).
\newcommand{\figuraApaisadaSinMarco}[5]{%
  \begin{figure}%
    \begin{center}%
    \includegraphics[angle=90,height=#1\textheight,#5]{#2}%
    \caption{#3}%
    \label{#4}%
    \end{center}%
  \end{figure}%
}
% Para las tablas
\newcommand{\otoprule}{\midrule [\heavyrulewidth]}
%
% Nuevo comando para tablas pequeñas (menos de una página).
\newcommand{\tablaSmall}[5]{%
 \begin{table}
  \begin{center}
   \rowcolors {2}{gray!35}{}
   \begin{tabular}{#2}
    \toprule
    #4
    \otoprule
    #5
    \bottomrule
   \end{tabular}
   \caption{#1}
   \label{tabla:#3}
  \end{center}
 \end{table}
}

%
%Para el float H de tablaSmallSinColores
\usepackage{float}

%
% Nuevo comando para tablas pequeñas (menos de una página).
\newcommand{\tablaSmallSinColores}[5]{%
 \begin{table}[H]
  \begin{center}
   \begin{tabular}{#2}
    \toprule
    #4
    \otoprule
    #5
    \bottomrule
   \end{tabular}
   \caption{#1}
   \label{tabla:#3}
  \end{center}
 \end{table}
}

\newcommand{\tablaApaisadaSmall}[5]{%
\begin{landscape}
  \begin{table}
   \begin{center}
    \rowcolors {2}{gray!35}{}
    \begin{tabular}{#2}
     \toprule
     #4
     \otoprule
     #5
     \bottomrule
    \end{tabular}
    \caption{#1}
    \label{tabla:#3}
   \end{center}
  \end{table}
\end{landscape}
}

%
% Nuevo comando para tablas grandes con cabecera y filas alternas coloreadas en gris.
\newcommand{\tabla}[6]{%
  \begin{center}
    \tablefirsthead{
      \toprule
      #5
      \otoprule
    }
    \tablehead{
      \multicolumn{#3}{l}{\small\sl continúa desde la página anterior}\\
      \toprule
      #5
      \otoprule
    }
    \tabletail{
      \hline
      \multicolumn{#3}{r}{\small\sl continúa en la página siguiente}\\
    }
    \tablelasttail{
      \hline
    }
    \bottomcaption{#1}
    \rowcolors {2}{gray!35}{}
    \begin{xtabular}{#2}
      #6
      \bottomrule
    \end{xtabular}
    \label{tabla:#4}
  \end{center}
}

%
% Nuevo comando para tablas grandes con cabecera.
\newcommand{\tablaSinColores}[6]{%
  \begin{center}
    \tablefirsthead{
      \toprule
      #5
      \otoprule
    }
    \tablehead{
      \multicolumn{#3}{l}{\small\sl continúa desde la página anterior}\\
      \toprule
      #5
      \otoprule
    }
    \tabletail{
      \hline
      \multicolumn{#3}{r}{\small\sl continúa en la página siguiente}\\
    }
    \tablelasttail{
      \hline
    }
    \bottomcaption{#1}
    \begin{xtabular}{#2}
      #6
      \bottomrule
    \end{xtabular}
    \label{tabla:#4}
  \end{center}
}

%
% Nuevo comando para tablas grandes sin cabecera.
\newcommand{\tablaSinCabecera}[5]{%
  \begin{center}
    \tablefirsthead{
      \toprule
    }
    \tablehead{
      \multicolumn{#3}{l}{\small\sl continúa desde la página anterior}\\
      \hline
    }
    \tabletail{
      \hline
      \multicolumn{#3}{r}{\small\sl continúa en la página siguiente}\\
    }
    \tablelasttail{
      \hline
    }
    \bottomcaption{#1}
  \begin{xtabular}{#2}
    #5
   \bottomrule
  \end{xtabular}
  \label{tabla:#4}
  \end{center}
}



\definecolor{cgoLight}{HTML}{EEEEEE}
\definecolor{cgoExtralight}{HTML}{FFFFFF}

%
% Nuevo comando para tablas grandes sin cabecera.
\newcommand{\tablaSinCabeceraConBandas}[5]{%
  \begin{center}
    \tablefirsthead{
      \toprule
    }
    \tablehead{
      \multicolumn{#3}{l}{\small\sl continúa desde la página anterior}\\
      \hline
    }
    \tabletail{
      \hline
      \multicolumn{#3}{r}{\small\sl continúa en la página siguiente}\\
    }
    \tablelasttail{
      \hline
    }
    \bottomcaption{#1}
    \rowcolors[]{1}{cgoExtralight}{cgoLight}

  \begin{xtabular}{#2}
    #5
   \bottomrule
  \end{xtabular}
  \label{tabla:#4}
  \end{center}
}




\graphicspath{ {./img/} }

% Capítulos
\chapterstyle{bianchi}
\newcommand{\capitulo}[2]{
	\setcounter{chapter}{#1}
	\setcounter{section}{0}
	\setcounter{figure}{0}
	\setcounter{table}{0}
	\chapter*{#2}
	\addcontentsline{toc}{chapter}{#2}
	\markboth{#2}{#2}
}

% Apéndices
\renewcommand{\appendixname}{Apéndice}
\renewcommand*\cftappendixname{\appendixname}

\newcommand{\apendice}[1]{
	%\renewcommand{\thechapter}{A}
	\chapter{#1}
}

\renewcommand*\cftappendixname{\appendixname\ }

% Formato de portada
\makeatletter
\usepackage{xcolor}
\newcommand{\tutor}[1]{\def\@tutor{#1}}
\newcommand{\course}[1]{\def\@course{#1}}
\definecolor{cpardoBox}{HTML}{E6E6FF}
\def\maketitle{
  \null
  \thispagestyle{empty}
  % Cabecera ----------------
\noindent\includegraphics[width=\textwidth]{cabecera}\vspace{1cm}%
  \vfill
  % Título proyecto y escudo informática ----------------
  \colorbox{cpardoBox}{%
    \begin{minipage}{.8\textwidth}
      \vspace{.5cm}\Large
      \begin{center}
      \textbf{TFG del Grado en Ingeniería Informática}\vspace{.6cm}\\
      \textbf{\LARGE\@title{}}
      \end{center}
      \vspace{.2cm}
    \end{minipage}

  }%
  \hfill\begin{minipage}{.20\textwidth}
    \includegraphics[width=\textwidth]{escudoInfor}
  \end{minipage}
  \vfill
  % Datos de alumno, curso y tutores ------------------
  \begin{center}%
  {%
    \noindent\LARGE
    Presentado por \@author{}\\ 
    en Universidad de Burgos --- \@date{}\\
        Tutores: \@tutor{} y\\
    Nuño Basurto Hornillos\\
  }%
  \end{center}%
  \null
  \cleardoublepage
  }
\makeatother


% Datos de portada
\title{GII-24.19 Contramedidas IoT mediante Reinforcement Learning \\Documentación Técnica}
\author{César Rodríguez Villagrá}
\tutor{Rubén Ruiz González}
\date{\today}

\begin{document}

\maketitle



\cleardoublepage



%%%%%%%%%%%%%%%%%%%%%%%%%%%%%%%%%%%%%%%%%%%%%%%%%%%%%%%%%%%%%%%%%%%%%%%%%%%%%%%%%%%%%%%%



\frontmatter


\clearpage

% Indices
\tableofcontents

\clearpage

\listoffigures

\clearpage

\listoftables

\clearpage

\mainmatter

\appendix

\apendice{Plan de Proyecto Software}

\section{Introducción}
En este apéndice se pretende mostrar el plan de proyecto software que se ha seguido durante el desarrollo del proyecto.

La implementación del plan de proyecto siguiendo una metodología ágil, ha permitido una mayor flexibilidad y adaptabilidad a los cambios que han surgido durante el desarrollo. En este caso, se ha utilizado principalmente la metodología de Kanban, debido a la flexibilidad que presenta en la organización y gestión del trabajo.

En cuanto a la estructura de este anexo, se explica la planificación temporal del proyecto y el estudio de viabilidad, tanto económica como legal. 

\subsection{Kanban}
La metodología Kanban~\cite{JuliaMartins2025} se ha utilizado para gestionar el flujo de trabajo del proyecto. En el que se usa el tablero de Kanban perdeterminado en los GitHub Projects, donde se han usado las siguientes columnas:
\begin{itemize}
    \item \textbf{Backlog:} En esta columna se encuentran todas las tareas que se han propuesto para el proyecto, pero que aún no se han comenzado a desarrollar, limitado a 5 tareas.
    \item \textbf{Ready:} En esta columna se encuentran todas las tareas que estaban en el \textit{Backlog}, y que están listas para ser desarrolladas.
    \item \textbf{In Progress:} En esta columna se encuentran todas las tareas que están en proceso de desarrollo, limitado a 3 tareas.
    \item \textbf{In Review:} En esta columna se encuentran todas las tareas que han sido completadas y están en revisión por parte de los autores del proyecto, limitado a 5 tareas.
    \item \textbf{Done:} En esta columna se encuentran todas las tareas que ya han sido completadas.
\end{itemize}

Todas las tareas que se han terminado quedan registradas en la columna \textit{Done}, y se han ido moviendo a lo largo de las columnas según se ha ido avanzando en el desarrollo del proyecto. En cada tarea, si forma parte de desarrollo de una parte de código o documentación, quedan registrado en la tarea una rama de desarrollo de GitHub, y una pull request asociada a la tarea en la que se han incorporado los cambios realizados a la rama principal. En cada pull request se han incluido comentarios creados pro GitHub Copilot, en la que se ha incluido una descripción de los cambios realizados, y una lista de los archivos que han sido modificados, y una revisión de errores o propuestas de mejora para implementar antes de la incorporación de cambios.

\section{Planificación temporal}
Cada semana, si es necesario se actualiza el tablero del proyecto.
Se han planificado cada semana una reunión de seguimiento del proyecto con los tutores, en la que se han revisado las tareas que se han realizado a lo largo de la semana. En la que se plantean propuestas de mejora y correcciones a las tareas realizadas.

Las reuniones semanales empezaron el día 6 de febrero de 2025 hasta el día 5 de junio de 2025, en el que cada jueves a las 10:00 se han realizado las reuniones, a excepción de festivos y vacaciones.

\subsection{Burn up}
El gráfico de Burn up muestra el progreso del proyecto a lo largo del tiempo, en el que se puede ver la cantidad de trabajo completado y el trabajo restante. En el eje horizontal se muestra el tiempo, y en el eje vertical se muestra la cantidad de trabajo completado. La línea morada representa el trabajo completado, y la línea verde representa el trabajo restante. El área entre las dos líneas representa el trabajo pendiente.
% #TODO Actualizar imágenes
\imagen{burn_up}{Gráfico de Burn up del proyecto}

\subsection{Diagrama de estado}
El diagrama de estado del proyecto muestra los diferentes estados en los que se encuentran actualmente el estado de las tareas del proyecto. Como el desarrollo del proyecto ya se ha terminado, el estado de las tareas se encuentra en "Done".
\imagen{project_status_chart}{Diagrama de estado del proyecto}

\subsection{Diagrama de trabajo por tema}
El diagrama de trabajo por tema muestra los diferentes temas en los que se ha trabajado a lo largo del proyecto, y la cantidad de trabajo realizado en cada uno de ellos. En el eje horizontal se muestran los diferentes temas, y en el eje vertical se muestra la cantidad de trabajo realizado.
Donde de izquierda a derecha se pueden ver: Lecturas de artículos,  Procesos de CI, Desarrollo de la web, con la api y el frontend, Documentación, Desarrollo y entrenamiento del agente.



\imagen{project_stimated}{Diagrama de trabajo por tema del proyecto}

El orden de los temas que han requerido más trabajo ha sido el siguiente:
\begin{enumerate}
    \item \textbf{Desarrollo y entrenamiento del agente \#7:} El tema que más trabajo ha requerido ha sido el desarrollo y entrenamiento del agente, ya que ha requerido una investigación previa, tanto para la implementación del agente y entorno como por el uso de aprendizaje por refuerzo.
    \item \textbf{Desarrollo de la web, con la api y el frontend \#6:} La parte de desarrollo de la aplicación web también ha requerido un gran esfuerzo, ya que se ha tenido que desarrollar tanto el backend como el frontend, así como la integración con la API REST y la base de datos.
    \item \textbf{Documentación \#68:}  Ya que se ha tenido que documentar tanto el código como el proyecto en general, así como la memoria y anexos del proyecto.
    \item \textbf{Procesos de CI \#47:} Los procesos de integración continua y calidad de código han sido claves en el inicio del proyecto para asegurar la calidad del código y la correcta ejecución de las pruebas, así como para facilitar el despliegue de la aplicación web.
    \item \textbf{Lecturas de artículos \#10:} La lectura de artículos ha sido necesaria para la investigación previa al desarrollo del proyecto, así como para la comprensión de los conceptos y técnicas utilizados en el proyecto, y la mejora de la función de recompensa del agente.
\end{enumerate}

\section{Estudio de viabilidad}
Realizar un estudio de viabilidad es una parte importante del futuro de un proyecto, ya que permite evaluar si el proyecto es viable desde diferentes perspectivas. En este caso, se ha realizado un estudio de viabilidad económica y legal, en el que se pretende estudiar cúanto dinero hace falta para realizar las tareas del proyecto, y si se cumplen las normativas legales necesarias para su desarrollo y despliegue.

\subsection{Viabilidad económica}

Los costes se dividen principalmente en dos categorías: costes de personal y costes de infraestructura.

Los costes de personal incluyen los salarios de los integrantes del proyecto. En este caso, se ha considerado ún unico empleado, en el cual ha trabajado en el proyecto de forma activa durante más de 5 menes, con una dedicación aproximada de 17 horas semanales, en el que dan unas horas totales de aproxiamdamente 380 horas de trabajo totales.

En resumen, la viabilidad económica del proyecto se ha evaluado teniendo en cuenta los costes de personal y de infraestructura, determinando que el proyecto es viable económicamente, ya que los costes totales son asumibles por la empresa o entidad que lo financie.

\subsection{Viabilidad legal}
Para la viabilidad legal del proyecto, se ha tenido en cuenta la normativa vigente en materia de protección de datos y propiedad intelectual. 
Las normativas aplicables más relevantes son:
\begin{itemize}
    \item \textbf{Reglamento General de Protección de Datos (RGPD):} Este reglamento establece las normas para la protección de datos personales en la Unión Europea.~\cite{RGPD}
    \item \textbf{Ley Orgánica de Protección de Datos y Garantía de los Derechos Digitales (LOPDGDD):} Esta ley complementa el RGPD en España y establece normas adicionales para la protección de datos personales.~\cite{LOPDGDD}	
\end{itemize}

Para analizar la viabilidad legal del proyecto, se debe tener en cuenta que el desarrollo y la implementación de un agente de aprendizaje por refuerzo encargado de la mitigación de ataques permitiendo o denegando el tráfico de red, debe garantizar que se cumplen las normativas vigentes.

En relación con la protección de datos personales del agente, cumple la normativa vigente, ya que no se recopilan ni procesan datos personales de los usuarios, ya que no obtiene información personal, ya que solo lee la capacidad de la cola, no accede a los datos de los paquetes. En la parte de la web, debido a que tampoco se recopilan datos personales de los usuarios, ya que en el registro de usuarios no se solicita información personal, y la única información que se almacena es el nombre de usuario y la contraseña, que se almacenan de forma segura en la base de datos, cumpliendo las mejores prácticas de seguridad, principalmente la autentificación y autorización de los usuarios, ya que utiliza Spring Security ~\cite{SpringSecurity}.
También usa servicios en la nube de terceros, como AWS~\cite{PPAWS}, Cloudflare~\cite{PPClouflare} y GitHub~\cite{PPGithub}, que cumplen con las normativas de protección de datos y seguridad, ya que se utilizan para el despliegue de la aplicación web, la gestión de dominios y seguridad del tráfico de la web y el control de versiones del código fuente, respectivamente.

En cuanto a las licencias de los paquetes y liberías usadas, se han utilizado principalmente:
\begin{itemize}
    \item \textbf{BSD-3-Clause:} Usada por más de 15 dependencias del proyecto.
    \item \textbf{BSD-2-Clause:} Usada por más de 13 dependencias del proyecto.
    \item \textbf{MIT:} Usada por 6 dependencias del proyecto.
    \item \textbf{Apache-2.0:} Usada por 5 dependencias del proyecto.
\end{itemize}
Todas estas licencias son compatibles con el proyecto, ya que permiten su uso y compercialización de la aplicación, siempre y cuando se cumplan las condiciones establecidas en cada licencia. En este caso, se ha cumplido con las condiciones de las licencias, ya que se ha incluido la información de las licencias en el proyecto, y no se ha modificado el código de las dependencias, ya que se han utilizado tal cual están en sus repositorios oficiales.


En resumen el proyecto cumple con la normativa vigente en materia de protección de datos y propiedad intelectual, permitiendo su desarrollo y despliegue sin problemas legales.
\apendice{Especificación de Requisitos}

\section{Introducción}

En este apéndice se presenta la especificación de requisitos del proyecto, que incluyen los objetivos generales, el catálogo de requisitos y la especificación de requisitos. La especificación de requisitos es una parte fundamental del desarrollo de software, ya que permite definir de forma clara y precisa las necesidades y expectativas del cliente.


\section{Objetivos generales}
La realización de este proyecto tiene como objetivo principal el desarrollo de un agente inteligente capaz de mitigar ataques de denegación de servicio en redes de computadoras, especialmente en redes de internet de las cosas (IoT). 
Para lograr este objetivo, se han establecido los siguientes objetivos:
\begin{itemize}
	\item Desarrollar un agente inteligente, y su entorno mediante aprendizaje por refuerzo que pueda mitigar ataques de denegación de servicio en redes.
	\item Hacer un estudio de los parámetros de la función de recompensa para optimizar el rendimiento del agente.
	\item Proporcionar una interfaz de usuario intuitiva para el usuario, que permita la visualización de estadísticas y gráficas relacionadas con el agente.
	\item Implementar un sistema que permita almacenar, consultar y eliminar modelos del agente.
\end{itemize}

\section{Catálogo de requisitos}
En el proyecto se han identificado los requisitos funcionales y no funcionales, desarrollados en la \textbf{Tabla~\ref{tabla:requisitosfuncionales}} y la \textbf{Tabla~\ref{tabla:requisitosnofuncionales}} respectivamente. Estos requisitos son esenciales para el correcto funcionamiento del sistema y deben ser cumplidos para garantizar los objetivos propuestos.

\tablaSmall{Tabla de requisitos funcionales}{l p{6em} p{15em} c}{requisitosfuncionales}{Número & Requerimiento & Descripción & Prioridad \\}{ 
RF1 & Registro de usuarios                  & Se permitirá registrarse nuevos usuarios al sistema.                                                                                                           & Alta  \\
RF2 & Inicio de sesión                      & Los usuarios registrados en el sistema podrán acceder con las credenciales que introdujo en el registro, pudiendo acceder a los modelos guardados previamente. & Alta  \\
RF3 & Panel de administración de usuarios   & Solo los usuarios con permisos de administrador podrán acceder a la gestión de usuarios, pudiendo borrar, y editar los usuarios registrados en el sistema.     & Media \\
RF4 & Importación y exportación de datos      & Se podrá importar y exportar los datos disponibles en la web.                                                                                                  & Baja  \\
RF5 & Documentación del sistema             & La web incluirá toda la información relacionada sobre el proyecto.                                                                                             & Media \\
RF6 & Visualización de los datos del agente & Se podrán visualizar datos relacionados con el agente, incluyendo datos de entrenamiento como del agente actuando.                                             & Alta  \\
} 

\tablaSmall{Tabla de requisitos no funcionales}{l p{6em} p{15em} c}{requisitosnofuncionales}{Número & Requerimiento & Descripción & Prioridad \\}{ 
RNF1 & Facilidad de uso                                      & Las interfaces de usuario deben de ser simples y adecuadas para todo tipo de usuarios.                                                          & Alta  \\
RNF2 & Rendimiento de la web                                 & El tiempo de carga de la web no debe de ser superior a 1 segundo para las operaciones más comunes.                                              & Media \\
RNF3 & Respuesta ante fallos                                 & Si ocurre un fallo, debe de tener el menor impacto posible en la web, mostrando un mensaje de error al usuario con el motivo del error.         & Alta  \\
RNF4 & Seguridad                                             & Un usuario solo debe de poder acceder solo a sus modelos y solo los usuarios con permisos de administrador podrán hacer la gestión de usuarios. & Alta  \\
RNF5 & Minimizar los paquetes que se descartan por el agente & El agente deberá de priorizar permitir el tráfico antes que denegarlo, intentando que los paquetes que se descartan sean los mínimos posibles.  & Alta  \\
} 

\section{Especificación de requisitos}
En esta sección se detallan los casos de uso del sistema, que describen las principales interacciones entre los usuarios y el sistema para cumplir con los requisitos funcionales identificados. Cada caso de uso incluye una descripción, precondiciones, acciones, postcondiciones, excepciones e importancia.

\begin{table}[p]
	\centering
	\begin{tabularx}{\linewidth}{ p{0.21\columnwidth} p{0.71\columnwidth} }
		\toprule
		\textbf{CU-1}    & \textbf{Registro de usuarios} \\
		\toprule
		\textbf{Versión}              & 1.0    \\
		\textbf{Autor}                & César Rodríguez Villagrá \\
		\textbf{Requisitos asociados} & RF1 \\
		\textbf{Descripción}          & Permite a un nuevo usuario crear una cuenta en el sistema proporcionando sus credenciales. \\
		\textbf{Precondición}         & 
		\begin{itemize}
			\item El usuario no debe estar registrado previamente.
			\item Disponibilidad del formulario de registro.
		\end{itemize} \\
		\textbf{Acciones}             & 
		\begin{enumerate}
			\item El usuario accede al formulario de registro.
			\item Introduce los datos requeridos.
			\item El sistema valida los datos.
			\item Si la validación es correcta, crea la cuenta y confirma el registro.
		\end{enumerate} \\
		\textbf{Postcondición}        & 
		\begin{itemize}
			\item El usuario queda registrado en el sistema.
			\item Puede iniciar sesión con sus credenciales.
		\end{itemize} \\
		\textbf{Excepciones}          & 
		\begin{itemize}
			\item Datos inválidos o incompletos: muestra mensaje de error.
			\item Nombre de usuario ya registrado: notificación para usar otro nombre.
		\end{itemize} \\
		\textbf{Importancia}          & Alta \\
		\bottomrule
	\end{tabularx}
	\caption{CU-1 Registro de usuarios.}
\end{table}


\begin{table}[p]
	\centering
	\begin{tabularx}{\linewidth}{ p{0.21\columnwidth} p{0.71\columnwidth} }
		\toprule
		\textbf{CU-2}    & \textbf{Inicio de sesión} \\
		\toprule
		\textbf{Versión}              & 1.0    \\
		\textbf{Autor}                & César Rodríguez Villagrá \\
		\textbf{Requisitos asociados} & RF2, RF1 \\
		\textbf{Descripción}          & Permite a usuarios registrados acceder al sistema introduciendo sus credenciales. \\
		\textbf{Precondición}         & 
		\begin{itemize}
			\item El usuario debe estar registrado.
			\item Disponibilidad del sistema.
		\end{itemize} \\
		\textbf{Acciones}             & 
		\begin{enumerate}
			\item El usuario abre el formulario de inicio de sesión.
			\item Introduce usuario y contraseña.
			\item El sistema valida las credenciales.
			\item Si son correctas, el usuario accede al sistema.
		\end{enumerate} \\
		\textbf{Postcondición}        & 
		\begin{itemize}
			\item El usuario está autenticado.
			\item Puede acceder a sus modelos.
		\end{itemize} \\
		\textbf{Excepciones}          & 
		\begin{itemize}
			\item Credenciales incorrectas: muestra mensaje de error.
		\end{itemize} \\
		\textbf{Importancia}          & Alta \\
		\bottomrule
	\end{tabularx}
	\caption{CU-2 Inicio de sesión.}
\end{table}


\begin{table}[p]
	\centering
	\begin{tabularx}{\linewidth}{ p{0.21\columnwidth} p{0.71\columnwidth} }
		\toprule
		\textbf{CU-3}    & \textbf{Gestión de usuarios} \\
		\toprule
		\textbf{Versión}              & 1.0    \\
		\textbf{Autor}                & César Rodríguez Villagrá \\
		\textbf{Requisitos asociados} & RF3 \\
		\textbf{Descripción}          & Los administradores pueden añadir, modificar o eliminar usuarios registrados en el sistema. \\
		\textbf{Precondición}         & 
		\begin{itemize}
			\item El usuario debe tener permisos de administrador.
			\item El sistema debe estar operativo.
		\end{itemize} \\
		\textbf{Acciones}             & 
		\begin{enumerate}
			\item El administrador accede al panel de gestión de usuarios.
			\item Puede consultar la lista de usuarios.
			\item Selecciona acciones como crear, editar o eliminar usuarios.
			\item El sistema aplica los cambios solicitados.
		\end{enumerate} \\
		\textbf{Postcondición}        & 
		\begin{itemize}
			\item La base de datos de usuarios queda actualizada.
			\item Cambios reflejados en futuras sesiones.
		\end{itemize} \\
		\textbf{Excepciones}          & 
		\begin{itemize}
			\item Intento de acción sin permisos: no permite el acceso.
			\item Error en la base de datos: muestra una notificación de fallo.
		\end{itemize} \\
		\textbf{Importancia}          & Media \\
		\bottomrule
	\end{tabularx}
	\caption{CU-3 Gestión de usuarios.}
\end{table}


\begin{table}[p]
	\centering
	\begin{tabularx}{\linewidth}{ p{0.21\columnwidth} p{0.71\columnwidth} }
		\toprule
		\textbf{CU-4}    & \textbf{Importación/Exportación de datos} \\
		\toprule
		\textbf{Versión}              & 1.0    \\
		\textbf{Autor}                & César Rodríguez Villagrá \\
		\textbf{Requisitos asociados} & RF4 \\
		\textbf{Descripción}          & Permite a los usuarios importar y exportar datos relacionados con los modelos y configuraciones. \\
		\textbf{Precondición}         & 
		\begin{itemize}
			\item Usuario autenticado.
			\item Formatos compatibles disponibles.
		\end{itemize} \\
		\textbf{Acciones}             & 
		\begin{enumerate}
			\item Si el usuario quiere descargar el modelo, selecciona el modelo deseado.
			\item El sistema procesa la operación. Utilizando el nombre adecuado.
			\item Muestra resultado y confirma la acción.
		\end{enumerate} \\
		\textbf{Postcondición}        & 
		\begin{itemize}
			\item Datos importados o exportados correctamente.
		\end{itemize} \\
		\textbf{Excepciones}          & 
		\begin{itemize}
			\item Archivo incompatible.
			\item Error durante la operación.
		\end{itemize} \\
		\textbf{Importancia}          & Baja \\
		\bottomrule
	\end{tabularx}
	\caption{CU-4 Importación/Exportación de datos.}
\end{table}


\begin{table}[p]
	\centering
	\begin{tabularx}{\linewidth}{ p{0.21\columnwidth} p{0.71\columnwidth} }
		\toprule
		\textbf{CU-5}    & \textbf{Visualización de datos y documentación} \\
		\toprule
		\textbf{Versión}              & 1.0    \\
		\textbf{Autor}                & César Rodríguez Villagrá \\
		\textbf{Requisitos asociados} & RNF1, RNF3 \\
		\textbf{Descripción}          & Permite a los usuarios consultar la documentación y datos relacionados con el agente RL a través de la web. \\
		\textbf{Precondición}         & 
		\begin{itemize}
			\item Usuario autenticado.
			\item Disponibilidad de la web.
		\end{itemize} \\
		\textbf{Acciones}             & 
		\begin{enumerate}
			\item El usuario accede al panel de visualización.
			\item Consulta documentación y estadísticas disponibles.
			\item Puede filtrar y buscar información relevante.
		\end{enumerate} \\
		\textbf{Postcondición}        & 
		\begin{itemize}
			\item Usuario obtiene la información solicitada.
		\end{itemize} \\
		\textbf{Excepciones}          & 
		\begin{itemize}
			\item Fallo en la carga de datos.
			\item Error de conexión.
		\end{itemize} \\
		\textbf{Importancia}          & Alta \\
		\bottomrule
	\end{tabularx}
	\caption{CU-5 Visualización de datos y documentación.}
\end{table}

\apendice{Especificación de diseño}

\section{Introducción}
La especificación de diseño es una parte fundamental del desarrollo de software, ya que permite definir de forma detallada el comportamiento y la estructura del sistema a desarrollar. En este apéndice se presenta la especificación de diseño del proyecto, que incluye el diseño de datos, el diseño arquitectónico y el diseño procedimental.

\section{Diseño de datos}
Para la transferencia de datos entre el frontend y el backend se han utilizado las clases de transferencia de datos (DTOs) con ayuda de Thymeleaf. En cambio, para la transferencia de datos entre el backend y la API se ha utilizado JSON, ya que es un formato ligero y sencillo para la comunicación entre sistemas.

La motivación de usar DTOs es:
\begin{itemize}
    \item \textbf{Reducción del acoplamiento:} Los DTOs permiten reducir el acoplamiento entre el frontend y el backend, lo que facilita la evolución y mantenimiento de la aplicación.
    \item \textbf{Seguridad y optimización de la transferencia de datos:} Los DTOs permiten transferir solo los datos necesarios entre el frontend y el backend, lo que reduce el tamaño de las peticiones y respuestas, además de proporcionar una capa adicional de seguridad, ya que por ejemplo en el DTO de los usuarios no incluirá la contraseña de los mismos.
    \item \textbf{Uso como formularios:} Los DTOs permiten definir de forma más clara los distintos tipos de validaciones de datos, por ejemplo, en el caso de los formularios, se pueden definir las validaciones necesarias para cada campo del formulario, útil por ejemplo para los formularios de registro.
\end{itemize}





Para la gestión de los datos de la base de datos se ha utilizado una base de datos relacional, en este caso MySQL. Se ha diseñado un esquema de base de datos que incluye las siguientes tablas:

\begin{itemize}
    \item \textbf{User}: Almacena información sobre los usuarios del sistema.
    \item \textbf{Role}: Define los roles de los usuarios.
    \item \textbf{TrainedModel}: Almacena los modelos de IA entrenados y datos relacionados.
\end{itemize}
Que se relacionan entre sí de la siguiente manera: \textbf{Figura~\ref{fig:dbrelaciones}}
\imagen{dbrelaciones}{Esquema de la base de datos del proyecto}

Para el intercambio de datos se han definido las siguientes clases de datos:
\imagen{dtos}{Clases de datos del proyecto}

En el que se muestran las clases, con sus constructores, atributos, métodos y propiedades, junto con su visibilidad y tipos.

\section{Diseño arquitectónico}

El diagrama de clases de alto nivel del proyecto muestra la estructura general del sistema, incluyendo las principales clases y sus relaciones. Este diagrama es una representación visual de cómo se organizan las clases en el sistema y cómo interactúan entre sí.
\imagen{relacionesClases}{Diagrama de clases del proyecto}

Para facilitar la visualización de las dependencias, se ha utilizado una matriz de dependencias agrupadas por paquetes. En el que los paquetes que se muestran en las primera columna, 
En cada columna, los paquetes siguen el mismo orden que en la primera columna, y se marcan en cada celda cúantas dependencias tiene un paquete con otro.

\imagen{matrizDep}{Matriz de dependencias del proyecto}
Como se puede observar en la matriz de la \textbf{Figura~\ref{fig:matrizDep}}, el paquete es.cesar.app.service tiene muchas dependencias, por lo que solo en este paquetes se han incluido las depencias de las clases individuales. En este caso la mayor cantidad de dependencias corresponden con la columna del paquete de los controladores(tercera columna). Por ejemplo, la clase UserService tiene dependencias con los paquetes de los controladores (es.cesar.app.controller) y configuración (es.cesar.app.config), es decir el paquetes de configuración(4 referencias) y controladores(29 referencias) dependen de la clase UserService, pero no al revés.

En la parte de la aplicación web, se ha utilizado el patrón de diseño Modelo-Vista-Controlador (MVC) para separar la lógica de negocio, la presentación y el control de flujo de la aplicación. Este patrón permite una mayor modularidad, facilita el mantenimiento y mejora la evolución del sistema. 
El mapa de navegación es una representación visual de las diferentes vistas de la aplicación y cómo se relacionan entre sí. En este caso, el mapa de navegación muestra las diferentes páginas de la aplicación web y cómo se accede a ellas desde otras páginas.

\imagen{mapanavegacion}{Mapa de navegación de la aplicación web}
La \textbf{Figura~\ref{fig:mapanavegacion}} muestra la estructura de navegación de la aplicación web, incluyendo las diferentes vistas y cómo se conectan entre sí.

\section{Diseño procedimental}
El diseño procedimental de este proyecto se ha centrado en estructurar el código de forma modular, facilitando así el desarrollo, mantenimiento y la posible extensión del sistema. En el que, para conseguir esto, se han utilizado buenas prácticas de programación y patrones de diseño. 

\subsection{Módulos}
Se ha realizado la división en 3 módulos:
\begin{itemize}
    \item RL
    \item app
    \item docs
\end{itemize}
En el que cada módulo tiene el mismo nombre que en el repositorio.
En el módulo de RL se encuentra todo lo relacionado con el aprendizaje por refuerzo, en el que proporciona en la carpeta principal lo relacionado con el entrenamiento y visualización de gráficas del modelo, y el sub-módulo de api, que proporciona una API REST para que desde la web pueda obtener datos de este módulo.

En el módulo app proporciona todo lo necesario de la aplicación web hecha con Spring Boot, en el que está subdividida en módulos, separando la configuración, los controladores, los archivos de transferencia de datos, los servicios etc.

En el módulo docs, proporciona todo lo relacionado con la generación de la documentación del proyecto, en el que encuentra los archivos base, y los archivos compilados.

Los módulos anteriores se unen mediante docker-compose (menos el de la documentación), creando un contenedor para la API, y otro para la aplicación web, junto a estos, también necesita un contenedor de una base de datos MySQL para persistir los datos de la web, garantizando un diseño modular y escalable.


\subsection{Patrones de diseño aplicados}
Para el diseño procedimental es importante el uso de patrones de diseño que faciliten la reutilización del código y la separación de responsabilidades. En este proyecto se han utilizado los siguientes patrones de diseño:
\begin{itemize}
    \item \textbf{Método plantilla(template method): } Este patrón se ha utilizado para definir los generadores de paquetes de datos para el entrenamiento de los modelos de IA. Definiendo una vez la estructura del generador de paquetes, y sobrescribiendo en las subclases los pesos de los tamaños de los paquetes de datos.
    \item \textbf{Estrategia(strategy): } En los generadores de tráfico, este patrón permite que se puedan definir diferentes estrategias de generación de paquetes de datos, en este caso de tráfico benignos y DoS, que se pueden intercambiar fácilmente sin afectar al resto del sistema.
    \item \textbf{Máquina de estados(state machine): } Este patrón se ha utilizado para gestionar los diferentes estados del tráfico, en el que hay estado de ataque y normal, intercambiando en cada estado el generador de paquetes de datos correspondiente. Permitiendo que se puedan definir diferentes estados y modificar las transiciones entre ellos, facilitando la gestión del flujo de la aplicación.
    \item \textbf{Singleton: } Este patrón se ha utilizado para la clase de instanciación de la API, que asegura que solo haya una instancia de la clase en toda la aplicación.
\end{itemize}
El uso de estos patrones permite una mayor modularidad, reutilización del código y facilita el mantenimiento y evolución del sistema, asegurando una mayor calidad de código.
\apendice{Documentación técnica de programación}

\section{Introducción}
Este apéndice contiene la documentación técnica de programación del proyecto, incluyendo la estructura de los directorios, manual del programador, cómo ejecutar el proyecto y las pruebas realizadas sobre el sistema. Esta documentación está destinada a las personas que quieran contribuir al proyecto o entender su funcionamiento interno.

\section{Estructura de directorios}
\label{sec:estructura-directorios}

La estructura de directorios del proyecto que se encuentra en el repositorio de GitHub principalmente consta de los siguientes directorios:
\begin{itemize}
    \item \textbf{app:} Directorio correspondiente al proyecto en Java con Spring Boot 3, que contiene todo el código y lógica de la aplicación web.
    \item \textbf{docs:} Directorio que contiene la documentación general del proyecto, que es donde se encuentra tanto los archivos fuente de LaTeX, como los pdfs de la memoria y el anexo.
    \item \textbf{RL:} Directorio que alberga el código fuente de la parte de Python, que contiene el código del agente de \textit{reinforcement learning} y la implementación sencilla de una API REST para la comunicación entre el agente y la aplicación web.
    \item \textbf{sql:} Que contiene los archivos sql con los que se iniciará la base de datos.
\end{itemize}
A parte de estas carpetas, en el directorio principal se puede encontrar los siguientes archivos:
\begin{itemize}
    \item \textbf{README.md:} Archivo que contiene una breve descripción del proyecto y las insignias (badges) de la calidad del código, la cobertura y el estado de las \hyperref[subsec:acciones]{acciones de GitHub}.    
    \item \textbf{.gitignore:} Archivo de configuración de Git, que indica los archivos y carpetas deben ser ignorados por el control de versiones.
    \item \textbf{sonar-project.properties:} Archivo de configuración de \hyperref[subsec:sonarqube]{SonarQube}, que contiene la configuración del análisis de calidad del código.
    \item \textbf{docker-compose.yml:} Archivo de configuración de Docker Compose, que permite levantar todos los contenedores necesarios  de Docker para el correcto funcionamiento del proyecto.
    \item \textbf{.env:} Archivo que contiene las variables de entorno necesarias para la ejecución del proyecto. El contenido de este archivo se desarrolla en más detalle en la \hyperref[subsec:politicas]{subsección de Políticas de seguridad}.
    \item \textbf{package.json:} Que es necesario para la ejecución de test con Jest, para los archivos de JavaScript.
    \item \textbf{pre\_containers.sh:} Contiene unos comandos necesarios para la correcta ejecución de los contenedores en el servidor de AWS EC2.
\end{itemize}


\section{Manual del programador}
\label{sec:manual-programador}

En la parte de gestión del \href{https://github.com/CesarRodrigu/GII-24.19-contramedidas-IoT-mediante-reinforcement-learning}{repositorio GitHub}, se ha realizado mediante un \href{https://github.com/users/CesarRodrigu/projects/6}{proyecto de GitHub} en el que se ha realizado la gestión de tareas(issues) y la planificación de las mismas, en las que la mayoría se han convertido en ramas de desarrollo para añadir los cambios propuestos.

El código del proyecto está documentado con comentarios, en los que se indica el funcionamiento de la clase, método, variable, etc. En el caso de la aplicación web, cuenta con un Javadoc en la carpeta de docs en el que se puede ver la documentación interna del código en formato HTML.

\subsection{Acciones de GitHub}
\label{subsec:acciones}
Las acciones de GitHub (GitHub Actions) son una herramienta de integración continua y entrega continua (CI/CD) que permite automatizar tareas en el flujo de trabajo de desarrollo de software. En este proyecto se han utilizado para realizar pruebas automáticas del código, pruebas de calidad y la Compilación de la documentación de formato LaTeX a PDF. A continuación se describen las acciones implementadas:
\begin{itemize}
    \item \textbf{Compilación de la documentación:} Se ha creado una acción que compila la documentación del proyecto en formato LaTeX a PDF. Esta acción se ejecuta cada vez que se realiza un push a la rama principal del repositorio en el que se han cambiado los archivos fuente de la propia documentación.
    \item \textbf{SonarQube:} Para las pruebas de calidad de código, se ha implementado SonarQube, que permite analizar el código en busca de errores, vulnerabilidades y problemas de calidad, según las reglas definidas en el proyecto. Cada vez que se añaden cambios a una incorporación de cambios o a la rama principal, se ejecuta una acción que analiza el código, ejecuta los test relacionados para obtener su cobertura y genera un informe de calidad. Este informe se puede consultar en la página de \href{https://sonarcloud.io/project/overview?id=CesarRodrigu_GII-24.19-contramedidas-IoT-mediante-reinforcement-learning}{SonarCloud del proyecto} y también se comenta en la \hyperref[subsec:sonarqube]{subsección dedicada a SonarQube}.

    \item \textbf{CodeQL:} GitHub CodeQL es una herramienta de análisis de código estático que permite detectar vulnerabilidades y errores en el código fuente. Cuando detecta un error, sugiere una posible solución mediante Copilot. Esta acción se ejecuta por cada incorporación de cambios al repositorio.
    \item \textbf{Dependabot:} Dependabot es una herramienta de GitHub que ayuda a mantener las dependencias del proyecto actualizadas y seguras, en este caso solo está configurada para que actualice las dependencias si se ha detectado una vulnerabilidad en alguna de ellas. Dependabot crea automáticamente una solicitud de incorporación de cambios para corregirlas, sugiriendo la versión que corrige esa vulnerabilidad.
\end{itemize}

\subsection{Políticas de seguridad}
\label{subsec:politicas}

En el repositorio se ha añadido un conjunto de reglas(ruleset), en el su función es asegurar que todos los cambios se realicen de manera controlada. En la que se ha definido una restricción para que no se pueda realizar un borrado de código no autorizado, y que todos los cambios realizados en el código se realicen mediante una solicitud de incorporación de cambios(pull request), en la que automáticamente realiza una solicitud de revisión de código a copilot y se ejecutan las \hyperref[subsec:acciones]{GitHub Actions ya descritas}.

También hay una opción de poder reportar fallos de seguridad en el apartado de \href{https://github.com/CesarRodrigu/GII-24.19-contramedidas-IoT-mediante-reinforcement-learning/security}{\textit{security}} si se han detectado en el código para que se puedan corregir lo más rápidamente posible.
Los escaneos de seguridad que se realizan sobre el código del repositorio principalmente se enfocan en la detección de paquetes vulnerables, vulnerabilidades en código y de secretos (credenciales y claves de acceso que no deberían estar en el código fuente).
Cabe destacar que el repositorio tiene en la carpeta raíz un archivo .env que contiene credenciales de acceso a la aplicación y este archivo no se debería subir al control de versiones. Sin embargo, en este caso se ha subido para que el código funcione correctamente, y que sea más fácil para un tercero ejecutarlo sin problemas, pero en un entorno de producción deberían estar como secretos de GitHub o solo disponer el archivo en el entorno local. Siendo altamente recomendable cambiar las credenciales cuando se vaya a desplegar el código en producción.
Es importante destacar que el archivo .env no contiene el secreto de la clave de la API de Groq, para evitar que se pueda utilizar la API de Groq sin autorización, ya que esta clave es personal y no debería ser compartida públicamente. En el caso de que se quiera utilizar la funcionalidad del chat de ayuda de la web se debe acceder a la página web publicada en la nube accesible en \url{https://www.cesarrv.com}.

\subsection{SonarQube}
\label{subsec:sonarqube}
SonarQube es una herramienta de análisis de código estático que permite detectar errores, vulnerabilidades y problemas de calidad en el código fuente y proporcionar métricas de calidad de software. En el proyecto de \href{https://sonarcloud.io/project/overview?id=CesarRodrigu_GII-24.19-contramedidas-IoT-mediante-reinforcement-learning}{SonarQube} se pueden ver cómo se ha evolucionado en los problemas de calidad de código, en la cobertura y en las duplicaciones de código. Proporciona una forma muy visual de ver la calidad del código mediante insignias(badges) incluidas en el \href{https://github.com/CesarRodrigu/GII-24.19-contramedidas-IoT-mediante-reinforcement-learning/blob/main/README.md}{README} del repositorio.

SonarQube trabaja con un conjunto de reglas que se deben cumplir para asegurar un mínimo de calidad de código del proyecto definida por el equipo de desarrollo. Estas reglas se pueden personalizar según las necesidades del proyecto, y se pueden añadir nuevas reglas o modificar las existentes. En este proyecto se han utilizado las reglas por defecto de SonarQube, ya que son suficientes para garantizar una buena calidad del código y principalmente porque son las únicas disponibles en el plan gratuito.


En la \textbf{Figura~\ref{fig:sonarGraphs}} se pueden ver algunas de las gráficas que proporciona SonarQube, que en este caso, muestran la evolución de la calidad del código, con las métricas de cobertura de pruebas(primera gráfica) y los issues detectados sin resolver(segundo gráfico). 
\imagen{sonarGraphs}{Gráficas de SonarQube del proyecto}

\section{Compilación, instalación y ejecución del proyecto}
\label{sec:compilacion}
Para la compilación e instalación del proyecto se puede hacer de forma general mediante Docker, ya que se ha creado un archivo \textit{docker-compose.yml} que permite levantar todos los contenedores necesarios para el correcto funcionamiento del proyecto. Para ello, se debe tener instalado Docker y Docker Compose en el sistema.

En el caso de que se quiera ejecutar la parte del agente de aprendizaje por refuerzo, se deben  seguir los mismos pasos descritos en la documentación de usuario en el apartado de instalación, que resume la instalación de las dependencias del proyecto en un entorno virtual, y un IDE asociado como Visual Studio Code~\cite{VSCode} o Jupyter-Notebook.

Si se quisiese compilar la parte de la web se recomienda el uso del IDE IntelliJ IDEA~\cite{intel}, en el que requiere tener instalado una versión de Java 21~\cite{java} y una versión compatible de Maven~\cite{maven}. Para instalar las dependencias de Java se debe usar el siguiente comando:
\begin{verbatim}
mvn clean install -f app/pom.xml
\end{verbatim}
Con esto ya se podrá editar y compilar la aplicación web. Para la ejecución del mismo se recomienda usar las herramientas que proporciona el IDE usado, o Docker.

Para la ejecución de pruebas se recomienda usar la GitHub Action asociada para evitar tener que instalar más dependencias. Si se quisiesen ejecutar en local, según el lenguaje se debería:
\begin{itemize}
    \item \textbf{Python} Ejecutar el comando pytest, en el que se pase como argumento la carpeta que contiene el código python a probar.
    \item \textbf{Java} ejecutar el comando:
    \begin{verbatim} 
    mvn clean verify -f app/pom.xml
    \end{verbatim}
    Donde maven se encarga de ejecutar las pruebas disponibles.
    \item \textbf{Pruebas de Selenium} instalar la extensión web(o la configuración de las pruebas en el proyecto web) e importar el archivo WebTests.side ubicado dentro de la carpeta de tests de Spring.
    \item \textbf{JavaScript} Se necesita tener instalado con yarn el paquete de Jest. y ejecutarlo con yarn jest.
\end{itemize}

\section{Pruebas del sistema}
\label{sec:pruebas}

Los distintos tipos de pruebas que se han realizado en el proyecto son:
\begin{itemize}
    \item \textbf{Pruebas unitarias:} Se han realizado pruebas unitarias tanto  en el código de Java, Python y JavaScript, utilizando JUnit 5, Pytest y Jest respectivamente. Todas estas pruebas se ejecutan automáticamente cada vez que se realiza un push al repositorio, mediante las \hyperref[subsec:acciones]{GitHub Actions} descritas anteriormente, que envían los resultados a Sonar, pudiendo consultar la cobertura de estas en la \hyperref[subsec:sonarqube]{página de SonarQube}. 
    En el caso de las pruebas unitarias de Java, se han realizado pruebas de integración con la base de datos, utilizando H2 como base de datos en memoria para evitar tener que realizar una conexión a una base de datos real. En el caso del agente de \textit{reinforcement learning}, se han realizado pruebas unitarias para comprobar que los componentes del agente y entorno funcionan correctamente.
    \item \textbf{Pruebas funcionales:} Se han realizado pruebas funcionales a través de Selenium, que permite, a traves de Selenium IDE, grabar pruebas simulando la interacción del usuario con la web. Gracias a la función de exportación de Selenium IDE las pruebas se pueden ejecutar tanto cargando el archivo WebTests.side, como ejecutando las pruebas de Java de la carpeta app/src/test/java/es/cesar/app/web (una vez ya iniciada la web) con un perfil de Spring Boot distinto a 'coverage', aceptando el perfil nativo. Se recomienda cargar el archivo WebTests.side en la extensión de Selenium IDE del navegador, ya que es más fácil de ejecutar y permite ver los resultados de las pruebas de forma más visual.
\end{itemize}

Para el desarrollo de las pruebas se ha utilizado la metodología \textbf{Arrange Act y Assert}, que en español sería Organizar, Actuar y Afirmar, que pretende seguir una misma estructura en todas las pruebas de primero declarar las variables necesarias, luego llamar a las funciones que se quieren probar y por último comprobar si los resultados son los correctos.

El objetivo de las pruebas es garantizar que la aplicación funciona correctamente y cumple con los requisitos funcionales y no funcionales definidos en la especificación del proyecto.
La integración continua de estas pruebas permite detectar y corregir errores cada vez que se incorporan cambios al código, lo que mejora la calidad del software y reduce el riesgo de introducir errores en el sistema.
Una alta cobertura de pruebas unitarias, como es en este caso, (superior al 85\% ) proporciona una garantía de calidad, significando que la mayoría del código se prueba de forma automática, facilitando la evolución del software.
\apendice{Documentación de usuario}

\section{Introducción}

\section{Requisitos de usuarios}

\section{Instalación}

\section{Manual del usuario}



\apendice{Anexo de sostenibilización curricular}

\section{Introducción}
Este anexo incluirá una reflexión personal del alumnado sobre los aspectos de la sostenibilidad que se abordan en el trabajo.
Se pueden incluir tantas subsecciones como sean necesarias con la intención de explicar las competencias de sostenibilidad adquiridas durante el alumnado y aplicadas al Trabajo de Fin de Grado.

Más información en el documento de la CRUE \url{https://www.crue.org/wp-content/uploads/2020/02/Directrices_Sosteniblidad_Crue2012.pdf}.

Este anexo tendrá una extensión comprendida entre 600 y 800 palabras.



\bibliographystyle{plain}
\bibliography{bibliografiaAnexos}

\end{document}
